% Generated draft: requires mathematical and source review.
\documentclass[11pt]{article}
\usepackage[margin=25mm]{geometry}
\usepackage{fontspec}
\setmainfont{Latin Modern Roman}
\usepackage{amsmath,amssymb,mathtools,unicode-math}
\setmathfont{Latin Modern Math}
\usepackage{enumitem,needspace,xurl,hyperref}
\setlength{\parindent}{0pt}
\setlength{\parskip}{5pt}
\setlength{\emergencystretch}{3em}
\newcommand{\N}{\mathbb{N}}
\newcommand{\Z}{\mathbb{Z}}
\newcommand{\Q}{\mathbb{Q}}
\newcommand{\R}{\mathbb{R}}
\newcommand{\C}{\mathbb{C}}
\newcommand{\F}{\mathbb{F}}
\newcommand{\A}{\mathbb{A}}
\newcommand{\PP}{\mathbb{P}}
\newcommand{\E}{\mathbb{E}}
\newcommand{\eps}{\varepsilon}
\newcommand{\dd}{\,\mathrm{d}}
\begin{document}
\section*{UnsolvedMath ID: 30002623}
\textbf{Source problem code:} catalogue entry
\subsection{Definitions and mathematical statement}
The supplied question is: \emph{Is every normal mod-2 current equivalent to a normal integral current?}

The record does not specify the ambient space, the dimension of the currents, the underlying theory of currents, or the meaning of equivalence modulo \(2\). Consequently, a fully quantified statement requires additional source clarification. The following is a conditional formulation of the representative question, rather than a resolution of those omissions.

Fix an ambient space \(X\), a dimension \(k\), and a theory of currents admitting reduction modulo \(2\). Under the finite-mass convention, a current \(T\) is normal when
\[
\mathbf M(T)+\mathbf M(\partial T)<\infty,
\]
where \(\partial T\) is its boundary and \(\mathbf M\) denotes mass. In the classical integral-current setting, an integral current is an integer-multiplicity rectifiable current whose boundary is also integer-multiplicity rectifiable, with both masses finite; thus it is normal. Write \(\mathcal I_k(X)\) for these currents. Write \(\mathcal N_k^{(2)}(X)\) for the mod-\(2\) currents satisfying
\[
\mathbf M_2(A)+\mathbf M_2(\partial A)<\infty,
\]
where \(\mathbf M_2\), the boundary operator, and the mod-\(2\) equivalence relation must be defined in the chosen theory. Denote reduction modulo \(2\) by \(\rho_2\).

With these conventions, the intended question has the form
\[
\forall A\in\mathcal N_k^{(2)}(X)\quad
\exists T\in\mathcal I_k(X)\quad \rho_2(T)=A.
\]
The source does not determine the range of \(X\) and \(k\) over which this assertion is intended, whether normality means finite mass or locally finite mass, or whether support restrictions apply. It also does not specify whether equivalence uses an algebraic quotient or an equivalence relation arising from a flat completion. These choices remain unresolved and are not supplied hypotheses of the original question.
\subsection{Short English statement}
Does every normal current with coefficients modulo \(2\) admit a normal integral representative in its modulo-\(2\) equivalence class? The ambient setting, dimensional scope, finiteness conventions, and equivalence relation require clarification. The supplied record labels the problem open; that status has not been independently verified.
\subsection{Sources}
\begin{itemize}
\item Local problem registry: \texttt{data/unsolvedmath/problems.json}, permanent UnsolvedMath ID \(30002623\), source problem code \texttt{catalogue entry}, title \emph{Integral Representatives of Normal Mod-Two Currents}. The supplied record gives no individual proposer and records the proposed year as \(2014\).
\item Supplied report citation: \emph{Calculus of Variations (2015)}, Oberwolfach Reports, pp.\ 1801--1860, DOI \texttt{10.4171/owr/2014/33}. Original Oberwolfach report: \url{https://doi.org/10.4171/owr/2014/33}. The supplied background describes Oberwolfach Reports as documenting mathematical workshops and their open-problem discussions. The citation's title year \(2015\), DOI year component \(2014\), and proposed year \(2014\) are retained without reconciliation.
\item Supplied related reference, labeled ``related quantitative nonorientability work'': \url{https://par.nsf.gov/servlets/purl/10054904}. No author, title, or precise supporting result is supplied for this reference.
\item Literature assessment dated \texttt{2026-08-21}: the supplied curation labels the exact representative question unresolved and states that related nonorientability and flat-chain results do not establish the stronger assertion. This is supplied dated curation derived from the OWR record, not an independently verified literature assessment. The record expressly calls for independent verification of time-sensitive status and cited claims.
\end{itemize}
\end{document}
